\documentclass[10pt,a4paper]{amsart}
\usepackage[margin=19mm]{geometry}
\usepackage{amsmath,amssymb,mathtools}
\usepackage[scr=rsfs,cal=euler]{mathalfa}
\usepackage{xfrac}
\usepackage[unicode,hidelinks,bookmarks=false]{hyperref}
\newcommand{\ie}{i.e.\ }
\newcommand{\cf}{cf.\ }
\newcommand{\resp}{respectively\ }
\newcommand{\Subsection}{\subsection}
\providecommand{\nobreakdash}{\nobreak\hbox{-}}
\setlength{\emergencystretch}{2em}
\multlinegap0pt
\usepackage{bm}
\usepackage{bbm}
\usepackage{dsfont}
\usepackage{xspace}
\usepackage{cancel}
\usepackage{mathtools}
\usepackage{amsthm}
\makeatletter
\@ifundefined{theorem}{\newtheorem{theorem}{Theorem}}{}
\@ifundefined{lemma}{\newtheorem{lemma}[theorem]{Lemma}}{}
\makeatother
\renewcommand{\tilde}{\widetilde}
\title[Perturbative Approach to Nonlinear Capacitance Matrix Formul]{Perturbative Approach to Nonlinear Capacitance Matrix Formulations}
\author{Habib Ammari, Clemens Thalhammer}
\date{Source: arXiv:2606.20821v1 (CC BY 4.0)}
\begin{document}
\maketitle

\noindent\emph{Source and licence.} Perturbative Approach to Nonlinear Capacitance Matrix Formulations. Version arXiv:2606.20821v1 (CC BY 4.0), distributed under the Creative Commons Attribution 4.0 International licence, as recorded in the arXiv licence metadata for this exact version and shown on the abstract page https://arxiv.org/abs/2606.20821v1. Full rights notice: licenses/arxiv-2606.20821-CC-BY-4.0.txt.\par\smallskip

\noindent\textit{Editorial scope.} Counted unit: the Lemma whose source label is \texttt{lem:approximate\_discrete} in the article (source0.tex lines 2120--2137, source label \texttt{lem:approximate\_discrete}), delivered together with the article's own complete original proof of it (source0.tex lines 2139--2273, one \texttt{proof} environment of 135 lines). No mathematics is written, repaired or re-derived: statement and proof are the article's own text line for line. Required context retained uncounted: eq:nonlinear\_problem (lines 570--584); eq:weak\_formulation (lines 633--635); eq:discrete\_nonlinear (lines 788--790); thm:converse (lines 1001--1026); lem:gradientbound (lines 2081--2089). Every definition, display and label of those lines is retained unchanged. Omitted from this package and not redistributed: 1--569, 585--632, 636--787, 791--1000, 1027--2080, 2090--2119, 2138--2138, 2274--3201. Nothing inside a retained excerpt is omitted or altered: every retained body is delivered exactly as the article's lines are written, so no omission receipt of any kind accompanies this package. Citations: the retained excerpts carry no citation command at all, so no bibliography entry of the article is transcribed and no reference is redistributed. Reference adaptation: none was needed and none was made; every reference inside the delivered unit resolves inside it, so no unresolved reference or citation is produced. Printed slips kept literal: no printed word, symbol, hypothesis or number of the counted statement or of its proof was altered. Layout and class: the article's own class is \texttt{amsart} with packages including bbm, geometry, inputenc, babel, fontenc, lmodern, amsmath, amssymb, amsthm, amsfonts, mathtools, mathdots, caption, subcaption; the wrapper is the lane's pinned \texttt{amsart} editorial wrapper at 10pt on A4, which restates the theorem-like environments the retained text needs and adds the article's own macro definitions that the retained text needs (\texttt{\textbackslash tilde}), with extra wrapper packages (bm, bbm, dsfont, xspace, cancel, mathtools). The delivered numbering is the unit's own. Assets: no retained span contains a figure environment, a graphics-inclusion command, a PDF-page inclusion, a table or a TikZ picture; no image file is copied into this package, the package declares no asset dependency and no image file is redistributed with it. Licence: version arXiv:2606.20821v1 of this article is distributed under the Creative Commons Attribution 4.0 International licence, as recorded in the arXiv licence metadata for this exact version and shown on the abstract page https://arxiv.org/abs/2606.20821v1. A full rights notice is delivered at licenses/arxiv-2606.20821-CC-BY-4.0.txt. Characters: every retained excerpt is pure ASCII, so the delivered engine, XeLaTeX with its default Latin Modern text font, reports no missing character.
\begin{equation}\label{eq:nonlinear_problem}
    \begin{cases}
        \Delta u + \dfrac{\omega^2}{v^2}\, u = 0
        & \text{in } \mathbb{R}^3 \setminus \overline{D}, \\[0.8em]
        \Delta u + \dfrac{\omega^2}{v_i^2}\, u
        + \sigma \dfrac{\omega^2}{v_i^2}\, |u|^2 u = 0
        & \text{in } D_i, \quad i \in \mathcal{I}, \\[0.8em]
        u\big|_+ = u\big|_-
        & \text{on } \partial D_i, \quad i \in \mathcal{I}, \\[0.8em]
        \delta_i \dfrac{\partial u}{\partial \nu}\bigg|_+
        = \dfrac{\partial u}{\partial \nu}\bigg|_-
        & \text{on } \partial D_i, \quad i \in \mathcal{I}, \\[0.8em]
        u \text{ is outgoing,}
    \end{cases}
\end{equation}

\begin{equation}\label{eq:weak_formulation}
    \sum_{i\in \mathcal{I}}\int_{D_i}\nabla u \cdot\nabla \overline{w} - \frac{\omega^2}{v_i^2}\left(u\overline{w} + \sigma \lvert u\rvert^2 u \overline{w}\right)\,dx - \delta \langle \mathcal{T}^{\omega/v} u, w\rangle_{\partial D} = 0
\end{equation}

\begin{equation}\label{eq:discrete_nonlinear}
        (\mathcal{C} - \omega_1^2 \mathbf{M})\mathbf{a} - \sigma \omega_1^2 \mathbf{N}(\mathbf{a}) = 0,
\end{equation}

\begin{theorem}[Continuous-to-discrete reduction]\label{thm:converse}
    Let $\{(\omega^\delta, u^\delta)\}_{\delta > 0} \subset \mathbb{C} \times H^1(D)$
    be a family of nontrivial solutions to~\eqref{eq:nonlinear_problem}
    satisfying the uniform bounds
    \begin{equation}\label{eq:uniform_bounds}
        \omega^\delta \sim \delta^{1/2},
        \qquad
        \lVert u^\delta \rVert_{L^2(D)}+ \lVert u^\delta \rVert_{L^4(D)} = \mathcal{O}(1),
    \end{equation}
    as $\delta \to 0$. Write $\omega^\delta = \delta^{1/2}\omega_1^\delta + o(\delta^{1/2})$
    and let $u_0^\delta := \Pi_d \, u^\delta \in \mathcal{V}_d$ denote the
    projection onto the space of resonator-wise constant functions.

    Then, after passing to a subsequence, if necessary, there exist
    $\omega_1 \in \mathbb{C} \setminus$ and
    $\mathbf{a} = (a_1, \dots, a_N)^\top \in \mathbb{C}^N \setminus$
    such that
    \begin{equation}
        \omega_1^\delta \to \omega_1,
        \qquad
        u_0^\delta|_{D_i} \to a_i,
    \end{equation}
    as $\delta \to 0$, and the limit $(\mathbf{a}, \omega_1)$ solves the
    discrete nonlinear capacitance
    system~\eqref{eq:discrete_nonlinear}. Moreover, if $\omega_1^2\notin\sigma(\mathbf{M}^{-1}\mathcal{C})$, the limit $\mathbf{a}$ is nonzero.
\end{theorem}

\begin{lemma}\label{lem:gradientbound}
    Under the above assumptions, for all $\delta$ sufficiently small,
    \begin{equation}
        \lVert \nabla u^\delta \rVert_{L^2(D)}^2
        \leq \frac{C_\omega^2}{v_{\min}^2}
        \bigl( \lVert u^\delta \rVert_{L^2}^2
          + \lVert u^\delta \rVert_{L^4}^4 \bigr) \delta.
    \end{equation}
\end{lemma}

\begin{lemma}\label{lem:approximate_discrete}
    Let $u_0^\delta := \Pi_d u^\delta \in \mathcal{V}_d$ and
    $\omega_1^\delta := \delta^{-1/2}\omega^\delta$ be as in
    Theorem~\ref{thm:converse}. Then,
    $\mathbf{a}^\delta := (u_0^\delta|_{D_i})_{i \in \mathcal{I}}$
    satisfies
    \begin{equation}
        \bigl\lVert
          (\mathcal{C} - (\omega_1^\delta)^2 \mathbf{M}) \mathbf{a}^\delta
          - \sigma (\omega_1^\delta)^2 \mathbf{N}(\mathbf{a}^\delta)
        \bigr\rVert_{\ell^2}
        \leq C
        \bigl( \lVert u^\delta \rVert_{L^2}^2
          + \lVert u^\delta \rVert_{L^4}^4 \bigr)^{1/2}
        \delta^{1/2},
    \end{equation}
    for some $C > 0$ independent of $\delta$.
\end{lemma}

\begin{proof}
    It suffices to show that
    \begin{equation}\label{eq:approx_dual}
        \bigl\lvert \bigl\langle
          (\mathcal{C} - (\omega_1^\delta)^2 \mathbf{M}) \mathbf{a}^\delta
          + \sigma (\omega_1^\delta)^2 \mathbf{N}(\mathbf{a}^\delta),
          \mathbf{w}
        \bigr\rangle \bigr\rvert
        \leq C
        \bigl( \lVert u^\delta \rVert_{L^2}^2
          + \lVert u^\delta \rVert_{L^4}^4 \bigr)^{1/2}
        \delta^{1/2}
    \end{equation}
    for all $\mathbf{w} \in \ell^2(\mathcal{I})$ with
    $\lVert \mathbf{w} \rVert_{\ell^2} = 1$. We identify $\mathbf{w}$
    with the element $w \in \mathcal{V}_d$ defined by $w|_{D_i} = w_i$,
    and write $\tilde{u}^\delta := u^\delta - u_0^\delta$. By the
    Poincar\'e--Wirtinger inequality on each $D_i$ and
    Lemma~\ref{lem:gradientbound},
    \begin{equation}\label{eq:tilde_bound}
        \lVert \tilde{u}^\delta \rVert_{L^2(D)}
        \leq C_P \lVert \nabla u^\delta \rVert_{L^2(D)}
        \leq \frac{C_P C_\omega}{v_{\min}}
        \bigl( \lVert u^\delta \rVert_{L^2}^2
          + \lVert u^\delta \rVert_{L^4}^4 \bigr)^{1/2}
        \delta^{1/2}.
    \end{equation}
    Testing the weak formulation~\eqref{eq:weak_formulation} against $w$
    and using $\nabla w = 0$ gives
    \begin{equation}\label{eq:tested_weak}
        \sum_{i \in \mathcal{I}} \frac{|\omega^\delta|^2}{v_i^2}
        \int_{D_i}
        \bigl( u^\delta + \sigma |u^\delta|^2 u^\delta \bigr)
        \overline{w}_i \, dx
        + \delta \langle \mathcal{T}^{\omega^\delta/v} u^\delta,
          w \rangle_{\partial D} = 0.
    \end{equation}
    We analyse each term, replacing $u^\delta$ by $u_0^\delta$ and
    $\mathcal{T}^{\omega^\delta/v}$ by $\mathcal{T}_0$. Throughout, $C$
    denotes a constant independent of $\delta$ that may change from
    line to line.

    For the linear volume term, since $\tilde{u}^\delta$ has zero mean
    on each $D_i$ and $w_i$ is constant on $D_i$,
    \begin{equation}
        \sum_{i \in \mathcal{I}} \frac{|\omega^\delta|^2}{v_i^2}
        \int_{D_i} u^\delta \overline{w}_i \, dx
        = \delta(\omega_1^\delta)^2
          \langle \mathbf{M} \mathbf{a}^\delta, \mathbf{w} \rangle.
    \end{equation}

    For the cubic volume term, the estimate
    $\bigl||a|^2 a - |b|^2 b\bigr| \leq 3(|a|^2 + |b|^2)|a - b|$
    and H\"older's inequality give
    \begin{equation}
        \begin{aligned}
            &\biggl\lvert
            \sum_{i \in \mathcal{I}} \frac{|\omega^\delta|^2}{v_i^2}
            \int_{D_i}
            \bigl(|u^\delta|^2 u^\delta
              - |u_0^\delta|^2 u_0^\delta\bigr)
            \overline{w}_i \, dx \biggr\rvert \\
            &\qquad\leq \frac{C_\omega^2 \delta}{v_{\min}^2}
            \bigl(
              \lVert u^\delta \rVert_{L^4}^2
              + \lVert u_0^\delta \rVert_{L^4}^2
            \bigr)
            \lVert \tilde{u}^\delta \rVert_{L^4}
            \lVert w \rVert_{L^4} \\
            &\qquad\leq C \lVert u^\delta \rVert_{L^4}^2
            \bigl( \lVert u^\delta \rVert_{L^2}^2
              + \lVert u^\delta \rVert_{L^4}^4 \bigr)^{1/2}
            \delta^{3/2},
        \end{aligned}
    \end{equation}
    where we used $\lVert u_0^\delta \rVert_{L^4}
    \leq \lVert u^\delta \rVert_{L^4}$,
    $\lVert \tilde{u}^\delta \rVert_{L^4}
    \leq \kappa_S \lVert \tilde{u}^\delta \rVert_{H^1}$,
    and~\eqref{eq:tilde_bound}. The leading-order contribution is
    $\delta \sigma (\omega_1^\delta)^2
    \langle \mathbf{N}(\mathbf{a}^\delta), \mathbf{w} \rangle$.

    For the boundary term, we write
    \begin{equation}
        \delta \langle \mathcal{T}^{\omega^\delta/v} u^\delta,
        w \rangle_{\partial D}
        = \delta \langle \mathcal{T}_0 u_0^\delta,
          w \rangle_{\partial D}
        + \delta \langle \mathcal{T}_0 \tilde{u}^\delta,
          w \rangle_{\partial D}
        + \delta \langle (\mathcal{T}^{\omega^\delta/v} - \mathcal{T}_0)
          u^\delta, w \rangle_{\partial D}.
    \end{equation}
    The first term equals
    $-\delta \langle \mathcal{C} \mathbf{a}^\delta,
    \mathbf{w} \rangle$. The second is bounded by
    \begin{equation}
        \delta \lVert \mathcal{T}_0 \rVert
        \lVert \tilde{u}^\delta \rVert_{H^1}
        \lVert w \rVert_{H^1}
        \leq C
        \bigl( \lVert u^\delta \rVert_{L^2}^2
          + \lVert u^\delta \rVert_{L^4}^4 \bigr)^{1/2}
        \delta^{3/2},
    \end{equation}
    and the third by
    \begin{equation}
        \delta \lVert \mathcal{T}^{\omega^\delta/v} - \mathcal{T}_0 \rVert
        \lVert u^\delta \rVert_{H^1}
        \lVert w \rVert_{H^1}
        \leq C
        \bigl( \lVert u^\delta \rVert_{L^2}^2
          + \lVert u^\delta \rVert_{L^4}^4 \bigr)^{1/2}
        \delta^2.
    \end{equation}

    Substituting into~\eqref{eq:tested_weak} and collecting the
    leading-order terms,
    \begin{equation}
        \delta \bigl\langle
          -\mathcal{C} \mathbf{a}^\delta
          + (\omega_1^\delta)^2 \mathbf{M} \mathbf{a}^\delta
          + \sigma (\omega_1^\delta)^2
            \mathbf{N}(\mathbf{a}^\delta),
          \mathbf{w}
        \bigr\rangle
        = O\bigl(
          (\lVert u^\delta \rVert_{L^2}^2
            + \lVert u^\delta \rVert_{L^4}^4)^{1/2}
          \delta^{3/2}
        \bigr).
    \end{equation}
    Dividing by $\delta$ yields~\eqref{eq:approx_dual}.
\end{proof}

\end{document}
